\track{深水区：四条被盲测指出「只有断言」的机制}{盲测反馈：下列四条是词典最薄处，单独补写机制，不重复既有词条。}

\begin{word}{单极子模空间与 Hecke 修正的机制}{monopole moduli and Hecke modifications}{monopole-hecke}{\S6.2, p.~81--82; \S9.1--9.2, p.~132--148; \S10.1--10.2, p.~148--155}
\mathdef{固定紧黎曼面 $C$、点 $p_0$ 与 $G$-丛 $E$。$E$ 在 $p_0$ 的 Hecke 修正是一对 $(E',\beta)$：$E'$ 为另一 $G$-丛，$\beta:E|_{C\setminus p_0}\to E'|_{C\setminus p_0}$ 为丛同构，等价关系是 $\beta$ 可延拓过 $p_0$。$G=U(N)$ 时局部 $E'=\bigoplus_iL_i\otimes\mathcal O(p_0)^{m_i}$，修正的型为支配权 ${}^Lw=(m_1\ge\cdots\ge m_N)$。修正空间是这些「丛连同 $p_0$ 外同构」的模空间，不是丛的同构类空间。}
\physdef{在四维 $\mathcal N=4$ 扭化理论中，沿直线 $L=\mathbb{R}\times\{z_0\}$ 插入 't Hooft 算子：只规定 $L$ 附近的奇异性而不规定和乐；$G=U(1)$ 时 $F\sim\frac{\mathrm{i}m}{2}\star_3\mathrm{d}(1/r)$，$m$ 为磁荷，非阿贝尔时奇异性按权 ${}^Lw$ 逐分量给出。它是无序算子，其「本征态」是边界条件而非四维态。}
\mathex{$G=U(2)$、${}^Lw=(1,0)$ 时修正只依赖 $p_0$ 处截面 $s_1$ 的复缩放，修正空间为 $\mathbb{CP}^1$，紧而光滑。${}^Lw=(2,0)$ 时它是加权射影空间 $\mathbb{P}(1,1,2)$ 挖去一点，紧化边界上又出现型 $(1,1)$ 的修正（§9.2, (9.20)--(9.21)）。}
\physex{紧化到 $\mathbb{R}\times C$ 后，算子的插入点 $z_0$ 就是 $p_0$。越过 $z_0$ 时全纯丛只可能在 $p_0$ 处跳变，跳变恰是 Hecke 修正：$G=U(1)$ 时 $T(m)$ 把 $L_-\mapsto L_+\cong L_-\otimes\mathcal O(p_0)^m$，$m$ 由 $\int_{S^2}c_1$ 读出（§9.1, (9.10) 附近）。}
\bridge{对应规则：算子 $T({}^Lw,z_0)$ 对应型 ${}^Lw$ 在 $z_0$ 的修正；「本征态」对应三维流形 $W=I\times C$ 上带该奇异性的 Bogomolny 解的椭圆边界条件；矩阵元由修正空间的相交上同调给出。自由参数是插入点 $z_0$ 与权 ${}^Lw$（即 ${}^LG$ 的最高权 ${}^LR$）。原文把修正空间等同于参数化奇异 Bogomolny 解的空间 $\bar{\mathcal Y}({}^Lw;z_0)$，并指出它同构于 Schubert 簇（§10.5, p.~182）。}
\intuit{一句话：'t Hooft 算子是「把磁荷钉在一点」的指令，而指令允许的奇点形状就是单极子，所以算子本征态的集合与单极子模空间是同一堆数据。}
\pitfall{别把这里的单极子模空间当成 $\mathbb{R}^3$ 上荷 $k$ 的孤立单极子模空间（那有 $4k$ 维）：这里是带 't Hooft 奇点、带椭圆边界条件的三维 Bogomolny 解空间，维数由 ${}^Lw$ 决定。更关键的是，「修正空间 = 单极子模空间」并不是算出来的定理：一个解限制到 $y$ 为常数的切片就给一次修正，反过来一次修正可按 $\partial_yA_{\bar z}=-\mathrm{i}D_{\bar z}\varphi_0$ 唯一延拓，两侧定义直接重合；真正的定理只是解空间光滑、维数等于椭圆指标。}
\end{word}

\begin{word}{Nahm 极点与 ADHMN 构造}{Nahm poles and the ADHMN construction}{nahm-pole}{\S2.1, p.~6（仅 Nahm 定理）；\S9.1, p.~132；\S10.5, p.~182；Nahm 方程原文未写出}
\mathdef{取区间 $I\subset\mathbb{R}$ 与 $k\times k$ 反 Hermite 矩阵函数 $T_1,T_2,T_3$。Nahm 方程是 $\mathrm{d}T_i/\mathrm{d}s=\frac12\varepsilon_{ijk}[T_j,T_k]$；Nahm 边界条件要求在 $I$ 的两端 $T_i$ 有单极点 $T_i\sim\rho(\sigma_i)/(2(s-s_0))$，其留数构成 $\mathfrak{su}(2)$ 的 $k$ 维不可约表示（$k=1$ 时表示平凡，端点没有极点）。数据模去 $U(k)$ 作用 $T_i\mapsto R^{-1}T_iR$ 得 Nahm 数据模空间；Nahm 定理断言它与 $\mathbb{R}^3$ 上荷 $k$ 的 BPS 单极子模空间一一对应。}
\physdef{极点的位置在物理上就是 Nahm 区间的端点（由 $v$ 定出，不是单极子位置）；留数是磁荷数据：它把未破缺的 $\mathfrak{su}(2)$ 嵌入 $\mathfrak{u}(k)$，表示的维数正是荷 $k$。ADHMN 构造：由 Nahm 数据解一维 Dirac 方程 $D^\dagger\psi=0$，取 $L^2$ 零模 $\psi(s,x)$，则 $\varphi_0(x)=\int s\,\psi^\dagger\psi\,\mathrm{d}s$、$A_j(x)=\int\psi^\dagger\partial_j\psi\,\mathrm{d}s$ 给出单极子的 Higgs 场与规范场。}
\mathex{$k=1$：$T_i$ 为常数，三个常数给出单极子位置（整体相位另计）。$k=2$：$T_i$ 为 $2\times2$ 矩阵，端点留数是二维不可约表示 $\sigma_i/2$，对应二单极子位形。区间 $I$ 的长度由渐近 Higgs 场 $|\varphi_0|\to v$ 固定，$v$ 是理论参数而不是模。}
\physex{$G=SU(2)$、荷 $1$：Nahm 数据只有三个常数，ADHMN 还原出 't Hooft--Polyakov 单极子，位置由常数读出、在核心处 $|\varphi_0|\to0$。荷 $2$：留数取二维表示，ADHMN 给出轴上对称的二单极子位形，其模空间维数 $4k=8$。}
\bridge{对应规则：Nahm 数据 $\leftrightarrow$ 单极子的谱数据；极点 $\leftrightarrow$ 渐近 Higgs 场给出的端点（广义数据里则是 't Hooft 算子的插入点）；留数 $\leftrightarrow$ 磁荷与表示；ADHMN $\leftrightarrow$ 由谱数据重构时空场。原文只在 §2.1 引用「Nahm 定理」（十维最大超对称），§9--§10 实际用的是 Bogomolny 方程与 ADHM 构造，从未写出 Nahm 方程本身。}
\intuit{一句话：Nahm 方程是单极子的「频域」描述，极点与留数像频域里的极点，ADHMN 就是把频域数据变回时空场的逆变换。}
\pitfall{常见误读是「$k$ 个极点对应荷 $k$ 的单极子」：标准 $SU(2)$ Nahm 数据的极点在区间两端、个数恒为 $2$，而荷是留数表示的维数（$k=1$ 时甚至没有极点）；只有广义 Nahm 数据、或把单极子逐个分开的退化极限里，极点才逐个承载一单位荷，此时极点位置才给出单极子沿 Nahm 方向的坐标。}
\end{word}

\begin{word}{膜到 D-模的翻译器}{the brane-to-D-module functor}{brane-dmod}{\S11.1, p.~185; \S11.2, p.~196; \S11.3, p.~198}
\mathdef{设 $M$ 为复流形，$T^*M$ 带全纯辛形式 $\Omega$，$\omega=\operatorname{Im}\Omega$。A-膜是三元组 $(Y,L,\nabla)$：Lagrangian 子流形 $Y\subset T^*M$、Hermite 线丛 $L$、酉联络 $\nabla$，满足 $F+\omega|_Y=0$；A-膜范畴的 Hom 是开弦态空间。取 $Y=T^*M$（充满空间）得典范余各向同性 A-膜 $B_{cc}$：其曲率满足 $(\omega^{-1}F)^2=-1$。函子 $\mathcal A\mapsto\operatorname{Hom}(B_{cc},\mathcal A)$ 把 A-膜送到 $K_M^{1/2}$ 上微分算子层 $\mathcal D_{K^{1/2}}$ 的模。}
\physdef{$B_{cc}$ 是 GL 扭化理论紧化后 $M_H$ 上 A-模型里的一个具体 A-膜；$(B_{cc},B_{cc})$ 开弦是边界 BRST 算子组成的结合代数。世界面瞬子被 $\omega-\mathrm{i}F$ 的非退化性排除，故开弦代数局部是层。$K^{1/2}$ 的来源是时间反演：$\mathsf T$ 把开弦代数变成相反代数，而 $\mathcal D_L^{\mathrm{op}}\cong\mathcal D_{L^{-1}\otimes K_M}$，自相反要求 $L^2\cong K_M$，故 $L=K_M^{1/2}$。}
\mathex{$G=U(1)$：Hitchin 模空间 $M=\operatorname{Jac}(C)\times B$，$B=H^0(C,\Omega^1)$，Hitchin 纤维化就是投影，纤维 $F_p\cong\operatorname{Jac}(C)$。取 $B'$ 为 Hitchin 纤维上的 type F 膜，则 $\operatorname{Hom}(B_{cc},B')$ 是取值于 $K_M^{1/2}$ 的 $(0,\bullet)$-形式，BRST 微分即 $\bar\partial$，其全纯上同调给出 $K_M^{1/2}$ 的全纯截面层，也就是 $\mathcal D_{K^{1/2}}$ 的一个模。}
\physex{$G=SU(2)$、亏格 $g\ge2$：$M=M(SU(2),C)$ 复维数 $3g-3$，$K_M^{1/2}$ 存在。取 $B'$ 为支撑在零截面 $\{\varphi=0\}$ 上的 A-膜，$\operatorname{Hom}(B_{cc},B')$ 的费米零模是 $T M$，量子化后得 $M$ 上旋量 $K_M^{1/2}\otimes\bigoplus_j\Omega^{0,j}$，其上半纯层就是 $M$ 上的扭化 D-模；这正是几何 Langlands 的特征 D-模。}
\bridge{对应规则：$K$ 复结构中的 A-膜范畴 $\to$ $\mathcal D_{K_M^{1/2}}$-模范畴，对象 $\mathcal A\mapsto\operatorname{Hom}(B_{cc},\mathcal A)$，$B_{cc}\mapsto\mathcal D_{K^{1/2}}$ 本身。扭化参数 $\Psi$ 进入方式是把 $\mathcal D_{K^{1/2}}$ 换成 $\mathcal D_{K^{1/2}\otimes\mathcal L^\Psi}$（$\mathcal L$ 为 $Bun_G$ 上行列式线丛，§11.3）。平方根线丛的选择不唯一：换一个 $K_M^{1/2}$ 得到的等价彼此相差一个 $\mathbb{Z}_2$-gerbe 的扭化，每个平方根选择给该 gerbe 一个平凡化（Frenkel 2009 综述脚注；§6.1 用它破坏 $A_1$ 与 $A_2$ 的对称）。}
\intuit{一句话：$B_{cc}$ 像「全空间的真空」，其他 A-膜是它上面的波函数，而作用在这些波函数上的微分算子就是 $B_{cc}$ 与自身的弦；给每个 A-膜配上它的波函数层，就把它翻译成了 D-模。}
\pitfall{别把 $\operatorname{Hom}(B_{cc},\mathcal A)$ 当成普通向量空间：它是复形或层，只有取上同调才是向量空间，函子的等价性也只在导出范畴层面成立。也别以为 $K^{1/2}$ 是自旋结构的选取：它是「代数自相反」的纯代数条件 $L^2\cong K_M$，与 $M$ 是否自旋无关。最后，KW06 §11 只给物理纲要（原文用 argued 而非 proved），严格化要靠 Nadler--Zaslow 一类几何方法。}
\end{word}

\begin{word}{三种拓扑扭化}{the three topological twists}{three-twistings}{\S3.1, p.~17--21; \S3.2, p.~22; \S5.1, p.~65}
\mathdef{四维 $\mathcal N=4$ 的 R-对称群为 $\operatorname{Spin}(6)\cong SU(4)$。扭化（twisting）是取同态 $\varkappa:\operatorname{Spin}(4)\to\operatorname{Spin}(6)$，把 $SU(4)$ 的作用经 $\varkappa$ 拉回；$\varkappa$ 的三种共轭类给出三类不等价的扭化。判据是扭化后保住的标量超荷个数：一个（半扭化）、两个（Vafa--Witten 型，由味对称联系、无连续族）、两个独立且反对易（GL 扭化）。GL 情形 $Q=uQ_\ell+vQ_r$，$(u:v)\in\mathbb{CP}^1$，$t=v/u$ 是唯一本质参数（$t\sim-t$ 等价）。}
\physdef{半扭化（文献中亦称 Donaldson--Witten 型）只留一个标量超荷 $Q$，$Q^2=0$，给出 Donaldson 型不变量；Vafa--Witten 扭化保留两个标量超荷，用于检验 S-对偶（文献 [45]）；GL 扭化保留两个独立、反对易的标量超荷，故任意线性组合仍是超对称，于是出现 $\mathbb{CP}^1$ 族，这是本文主线。}
\mathex{GL 扭化下玻色场变成规范场 $A$、伴随 $1$-形式 $\varphi$ 和复标量 $\sigma,\bar\sigma$：$6_v=(2,2)^0\oplus(1,1)^2\oplus(1,1)^{-2}$（§3.1, (3.4)）。$t=\pm\mathrm{i}$ 时方程化为 $F(A+\mathrm{i}\varphi)=0$，即复平坦联络，模空间为 $Y_M(G)$；$t$ 实时为实椭圆方程族，以 $\mathbb{RP}^1$ 为参数（§3.2, (3.29)--(3.30)）。}
\physex{$t=\pm\mathrm{i}$：复结构 $\pm J$ 中的 B-模型。$t=\pm1$：$(w_+,w_-)=(\mp1,\pm1)$，故 $J_+=-J_-=\pm K$，得复结构 $K$ 中的 A-模型。$t=0$ 或 $\infty$：$(w_+,w_-)=(0,\infty)$，得复结构 $I$ 中的 A-模型（这三行由 §5.1 的判据与 (5.11) 直接读出）。S-对偶把 $t$ 映为 $(\tau/|\tau|)t$（§3.1, (3.13)）。}
\bridge{对应规则：三类扭化 $\leftrightarrow$ 保住的标量超荷数（一个、两个无族、两个成族）；GL 族参数 $t$ $\leftrightarrow$ $M_H$ 的复结构对 $(w_+,w_-)$ $\leftrightarrow$ A/B-模型；$t$ 与 S-对偶的 $\tau$ 按 (3.13) 联动。原文位置：§3.1 p.~17（三类扭化）、§3.2 p.~22（$t$ 族方程）、§5.1 p.~65（$w_\pm$ 与 A/B 判据）、§5.3（S-对偶把 $t=\mathrm{i}$ 送到 $t=1$）。}
\intuit{一句话：扭化是把 R-对称群的内部转动与时空转动绑在一起，使某个超荷变成标量；三种绑法里只有 GL 绑法让两个标量超荷可以任意混合，于是参数化出 $\mathbb{CP}^1$。}
\pitfall{最常见的伪问题是「三种扭化各给 $I,J,K$ 一个复结构」：$I,J,K$ 不是三种扭化的标签，而是 GL 扭化这一族在不同 $t$ 处读出的复结构（$t=\pm\mathrm{i}$ 给 $J$，$t=\pm1$ 给 $K$，$t=0,\infty$ 给 $I$），三种扭化中只有 GL 带连续族。第二个误解是以为 §5.1 用到 Vafa--Witten 扭化：§5.1 的 $J_\pm$ 族是 GL 扭化的二维约化。原文对前两类只说「与 Donaldson 理论类似」并点名文献 [45]，并未给出惯用名，Donaldson--Witten 型与 Vafa--Witten 型是按文献惯例补的。}
\end{word}
